\documentclass[12pt, a4paper]{report}

% Подключение необходимых пакетов для русского языка и математики
\usepackage[utf8]{inputenc}
\usepackage[T2A]{fontenc}
\usepackage[english, russian]{babel}
\usepackage{amsmath, amsfonts, amssymb} % Математические пакеты
\usepackage{graphicx} % Для вставки графиков
\usepackage{geometry} % Настройка полей
\usepackage{hyperref} % Гиперссылки
\usepackage{csquotes} % Оформление цитат
\usepackage[backend=biber,style=gost-numeric,sorting=none]{biblatex} % Библиография по ГОСТ
\usepackage{titlesec} % Настройка заголовков
\usepackage{fancyhdr} % Колонтитулы
\usepackage{setspace} % Межстрочный интервал

% Настройка полей (по ГОСТ для курсовых/дипломных работ)
\geometry{top=2cm, bottom=2cm, left=3cm, right=1.5cm}

% Настройка межстрочного интервала (полуторный)
\onehalfspacing

% Настройка заголовков
\titleformat{\section}
{\normalfont\fontsize{14}{16}\bfseries}
{\thesection}{1em}{}

\titleformat{\subsection}
{\normalfont\fontsize{13}{15}\bfseries}
{\thesubsection}{1em}{}

\titleformat{\subsubsection}
{\normalfont\fontsize{12}{14}\bfseries}
{\thesubsubsection}{1em}{}

% Настройка колонтитулов
\pagestyle{fancy}
\fancyhf{}
\fancyhead[R]{\thepage}
\fancyhead[L]{}
\renewcommand{\headrulewidth}{0pt}

\addbibresource{references.bib}

\begin{document}

% Титульный лист (без номера страницы)
\thispagestyle{empty}
\begin{center}
\vspace*{2cm}

{\Large \textbf{НАЗВАНИЕ УЧЕБНОГО ЗАВЕДЕНИЯ}}

\vspace{1cm}

{\large Кафедра [название кафедры]}

\vspace{3cm}

{\large \textbf{ПРОБЛЕМАТИКА РАБОТЫ И ТЕОРЕТИЧЕСКИЙ АНАЛИЗ}}

\vspace{1cm}

{\large [Название работы]}

\vspace{3cm}

\begin{flushright}
Выполнил: [ФИО студента] \\
Группа: [номер группы] \\
Научный руководитель: [ФИО руководителя]
\end{flushright}

\vspace{2cm}

{\large [Город, год]}

\end{center}

\newpage

% Содержание
\tableofcontents
\newpage

% Основной текст начинается с арабской нумерации
\setcounter{page}{3}

\section{Введение}

В фокусе данного исследования находится задача \textit{RPML (Repayment Policy for Multiple Loans)} — оптимизация погашения множественных долговых обязательств в условиях бюджетных ограничений. Актуальность задачи обусловлена низкой эффективностью популярных эвристических стратегий, широко используемых в персональных финансах, таких как \textit{Debt Avalanche} (погашение кредитов с максимальной процентной ставкой) и \textit{Debt Snowball} (погашение кредитов с минимальным балансом).

\section{Проблематика работы и теоретический анализ}

\subsection{Проблематика: ограниченность эвристических подходов}

Теоретический анализ показывает, что данные эвристики приводят к субоптимальным решениям из-за игнорирования нелинейных ограничений реальных банковских продуктов:
\begin{itemize}
    \item \textbf{Штрафные ставки (Penalty Rates):} Пропуск минимального платежа по одному кредиту ради досрочного погашения другого может активировать штрафную ставку (default rate), которая нивелирует экономию на процентах.
    \item \textbf{Кассовые разрывы (Cash Flow Gaps):} Детерминированные стратегии не учитывают волатильность доходов. Оптимальная стратегия часто требует контринтуитивного накопления «буферного фонда» (сбережений) вместо досрочного погашения, чтобы избежать дефолта в будущие периоды низкой ликвидности.
\end{itemize}

Согласно исследованиям, следование стратегии Debt Avalanche может приводить к переплате в среднем на 4.6\%, а в отдельных сценариях — до 40\% от общей суммы долга по сравнению с математическим оптимумом \cite{rios2017}.

\subsection{Теоретический анализ сложности}

Ключевым теоретическим выводом является принадлежность задачи RPML к классу \textbf{NP-трудных} (NP-hard) задач. 

Rios-Solis и соавторы (2017) доказали сложность задачи через полиномиальную сводимость (reduction) к известной задаче о многих рюкзаках (\textit{Multiple Knapsack Problem}). Формально, если $N$ — количество кредитов, а $T$ — временной горизонт планирования, то пространство поиска решений растет экспоненциально как $O(2^{N \times T})$. Это означает отсутствие полиномиального алгоритма, гарантирующего нахождение глобального оптимума за приемлемое время для произвольного числа кредитов, что обосновывает необходимость использования методов математического программирования.

\subsection{Обзор существующих решений (SOTA)}

Эволюция методов решения задачи RPML прошла несколько этапов, от упрощенных эвристик до сложных моделей смешанного целочисленного программирования (MILP).

\subsubsection{Ранние подходы и аппроксимации}
Первые попытки точного решения были предприняты Gupta et al., 1987, предложившим алгоритм ветвей и границ (\textit{Branch-and-Bound}) для упрощенных моделей без учета штрафных ставок \cite{gupta1987}. Позднее Нг (Ng et al., 2012) переформулировал задачу как проблему планирования на единичном станке с ухудшающимся временем обработки (\textit{scheduling with deteriorating processing times}), предложив полиномиальный алгоритм. Однако данная модель не учитывала обязательные минимальные платежи, что делало её неприменимой для кредитных карт \cite{ng2012}.

\subsubsection{Текущий Baseline. Детерминистическая MILP модель}
На текущий момент SOTA (State of the Art) решением является модель, предложенная Риос-Солисом (Rios-Solis et al., 2017). Авторы сформулировали задачу как проблему смешанного целочисленного линейного программирования (MILP), целью которой является минимизация общей суммы выплат $Z$:

\begin{equation}
    Z = \sum_{t=1}^T \sum_{j=1}^N X_{jt} \rightarrow \min,
    \label{eq:objective}
\end{equation}

где $X_{jt}$ — платеж по кредиту $j$ в месяц $t$. Модель учитывает:
\begin{itemize}
    \item Ограничения бюджета домохозяйства;
    \item Динамику баланса долга с учетом капитализации процентов;
    \item Бинарные переменные, активирующие штрафные ставки при пропуске минимального платежа.
\end{itemize}

Недавние исследования (Friedetzky et al., 2024) расширили данный подход на сетевые структуры (\textit{Interval Debt Model}), однако они сфокусированы на межбанковских расчетах, а не на персональных финансах \cite{friedetzky2024}.

\subsection{Недостатки SOTA и личный вклад в разработку}

Несмотря на математическую строгость, существующие SOTA решения (в частности, модель Rios-Solis) имеют критические ограничения для практического применения:
\begin{enumerate}
    \item \textbf{Детерминизм:} Базовая модель предполагает полное знание будущих доходов на весь горизонт планирования, игнорируя стохастическую природу заработка.
    \item \textbf{Риск-нейтральность:} Целевая функция минимизирует математическое ожидание затрат, игнорируя «хвостовые» риски (tail risks) банкротства в наихудших сценариях.
\end{enumerate}

\textbf{Личный вклад автора} заключается в устранении данных ограничений через разработку стохастической модификации MILP-модели. В отличие от детерминистического подхода, предлагается:
\begin{itemize}
    \item Введение сценарного анализа по методу Монте-Карло для моделирования неопределенности доходов.
    \item Интеграция метрики \textbf{CVaR (Conditional Value at Risk)} в целевую функцию для минимизации средних потерь в 5\% наихудших сценариев.
    \item Реализация механизма управления динамическим буферным фондом.
\end{itemize}

Рабочая гипотеза предполагает, что данный подход позволит снизить вероятность кассового разрыва (\textit{Cash Shortfall Rate}) на $\ge 45\%$ относительно детерминистического решения RPML.

\newpage

% Список литературы
\printbibliography[title={Список литературы}]

\end{document}